% GENERATED FILE. Edit canonical lessons, then run npm run generate:books.
% canonical-source-hash: fnv1a64:fbb43714cfd0996f
% canonical-lessons: GU-C234-vachardu, GU-C234-bhabharvu, GU-C234-dohvu, GU-C234-ghetu, GU-C234-carvu

\chapter{Milking time}
\label{ch:gu-milking-time}
\hlchaptermodality{234}

\begin{chapteropening}
\textbf{By the end of this chapter:} \emph{I can say who milked the cow and where the sheep grazed.}

Complete the last lesson of chapter 234: I can say who milked the cow and where the sheep grazed.
\end{chapteropening}

\section[\texorpdfstring{\gu{વાછરડું} (vācharḍũ)}{વાછરડું (vācharḍũ)}]{\gu{વાછરડું} (vācharḍũ) — calf}

\label{lesson:GU-C234-vachardu}



\begin{quote}
Before the new one: say the Gujarati for ox, then the Gujarati for to plough.
\end{quote}

\subsection*{You'll want to know: \gu{વાછરડું}}
\textbf{\gu{વાછરડું}} — \emph{vācharḍũ} — \textquotedblleft{}calf\textquotedblright{}.

\begin{quote}
\textbf{\gu{વાછરડું} \gu{ગાય} \gu{પાસે} \gu{ગયું}.} — \emph{vācharḍũ gāy pāse gayũ} — \textquotedblleft{}The calf went to the cow.\textquotedblright{}
\end{quote}

\begin{cousinweb}
\textbf{\gu{વાછરડું}} (\emph{vācharḍũ}) is a tadbhava. It comes from Sanskrit \emph{vatsatara}, \textquotedblleft{}young bull or goat before weaning\textquotedblright{}, built on \emph{vatsa}, \textquotedblleft{}calf\textquotedblright{} (CDIAL 11241).
\end{cousinweb}

\begin{grammarlens}[title={What you've built}]
One more word for everyday life at A2.
\end{grammarlens}

\subsection*{Your turn}
\begin{itemize}
\raggedright
  \item \emph{Say it:} \emph{vācharḍũ}
  \item \emph{Say it:} \emph{vācharḍũ}, once more
  \item \emph{Say it:} \emph{vācharḍũ gāy pāse gayũ} — The calf went to the cow.
  \item \emph{Recall:} say the Gujarati for ox, then the Gujarati for to plough, then say \emph{vācharḍũ} again
\end{itemize}

\begin{tcolorbox}[breakable,colback=teal!4,colframe=teal!35!black,title={Before you move on}]
What does \gu{વાછરડું} mean? (\textquotedblleft{}Calf\textquotedblright{}.) Say it once more.
\end{tcolorbox}
\section[\texorpdfstring{\gu{ભાંભરવું} (bhā̃bharvũ)}{ભાંભરવું (bhā̃bharvũ)}]{\gu{ભાંભરવું} (bhā̃bharvũ) — to moo, to low}

\label{lesson:GU-C234-bhabharvu}



\begin{quote}
Before the new one: say the Gujarati for to plough, then the Gujarati for calf.
\end{quote}

\subsection*{You'll want to know: \gu{ભાંભરવું}}
\textbf{\gu{ભાંભરવું}} — \emph{bhā̃bharvũ} — \textquotedblleft{}to moo, to low\textquotedblright{}.

\begin{quote}
\textbf{\gu{સવારે} \gu{વાછરડું} \gu{ભાંભર્યું}.} — \emph{savāre vācharḍũ bhā̃bharyũ} — \textquotedblleft{}The calf mooed in the morning.\textquotedblright{}
\end{quote}

\begin{cousinweb}
The origin of \textbf{\gu{ભાંભરવું}} (\emph{bhā̃bharvũ}) is not settled. It is probably imitative of the sound a cow makes, and CDIAL lists no Gujarati form like it.
\end{cousinweb}

\begin{grammarlens}[title={What you've built}]
One more word for everyday life at A2.
\end{grammarlens}

\subsection*{Your turn}
\begin{itemize}
\raggedright
  \item \emph{Say it:} \emph{bhā̃bharvũ}
  \item \emph{Say it:} \emph{bhā̃bharvũ}, once more
  \item \emph{Say it:} \emph{savāre vācharḍũ bhā̃bharyũ} — The calf mooed in the morning.
  \item \emph{Recall:} say the Gujarati for to plough, then the Gujarati for calf, then say \emph{bhā̃bharvũ} again
\end{itemize}

\begin{tcolorbox}[breakable,colback=teal!4,colframe=teal!35!black,title={Before you move on}]
What does \gu{ભાંભરવું} mean? (\textquotedblleft{}To moo, to low\textquotedblright{}.) Say it once more.
\end{tcolorbox}
\section[\texorpdfstring{\gu{દોહવું} (dohvũ)}{દોહવું (dohvũ)}]{\gu{દોહવું} (dohvũ) — to milk}

\label{lesson:GU-C234-dohvu}



\begin{quote}
Before the new one: say the Gujarati for calf, then the Gujarati for to moo.
\end{quote}

\subsection*{You'll want to know: \gu{દોહવું}}
\textbf{\gu{દોહવું}} — \emph{dohvũ} — \textquotedblleft{}to milk\textquotedblright{}.

\begin{quote}
\textbf{\gu{સવારે} \gu{મેં} \gu{ગાય} \gu{દોહી}.} — \emph{savāre mẽ gāy dohī} — \textquotedblleft{}I milked the cow in the morning.\textquotedblright{}
\end{quote}

\begin{cousinweb}
\textbf{\gu{દોહવું}} (\emph{dohvũ}) is a tadbhava. It comes from \emph{*dohati}, \textquotedblleft{}milks\textquotedblright{}, a form of the Sanskrit root \emph{duh}, \textquotedblleft{}milk\textquotedblright{} (CDIAL 6592).
\end{cousinweb}

\begin{grammarlens}[title={What you've built}]
One more word for everyday life at A2.
\end{grammarlens}

\subsection*{Your turn}
\begin{itemize}
\raggedright
  \item \emph{Say it:} \emph{dohvũ}
  \item \emph{Say it:} \emph{dohvũ}, once more
  \item \emph{Say it:} \emph{savāre mẽ gāy dohī} — I milked the cow in the morning.
  \item \emph{Recall:} say the Gujarati for calf, then the Gujarati for to moo, then say \emph{dohvũ} again
\end{itemize}

\begin{tcolorbox}[breakable,colback=teal!4,colframe=teal!35!black,title={Before you move on}]
What does \gu{દોહવું} mean? (\textquotedblleft{}To milk\textquotedblright{}.) Say it once more.
\end{tcolorbox}
\section[\texorpdfstring{\gu{ઘેટું} (gheṭũ)}{ઘેટું (gheṭũ)}]{\gu{ઘેટું} (gheṭũ) — sheep}

\label{lesson:GU-C234-ghetu}



\begin{quote}
Before the new one: say the Gujarati for to moo, then the Gujarati for to milk.
\end{quote}

\subsection*{You'll want to know: \gu{ઘેટું}}
\textbf{\gu{ઘેટું}} — \emph{gheṭũ} — \textquotedblleft{}sheep\textquotedblright{}.

\begin{quote}
\textbf{\gu{ઘેટું} \gu{ટેકરીથી} \gu{નીચે} \gu{આવ્યું}.} — \emph{gheṭũ ṭekrīthī nīce āvyũ} — \textquotedblleft{}The sheep came down from the hill.\textquotedblright{}
\end{quote}

\begin{cousinweb}
The origin of \textbf{\gu{ઘેટું}} (\emph{gheṭũ}) is not settled: CDIAL lists no Gujarati form like it among its words for sheep.
\end{cousinweb}

\begin{grammarlens}[title={What you've built}]
One more word for everyday life at A2.
\end{grammarlens}

\subsection*{Your turn}
\begin{itemize}
\raggedright
  \item \emph{Say it:} \emph{gheṭũ}
  \item \emph{Say it:} \emph{gheṭũ}, once more
  \item \emph{Say it:} \emph{gheṭũ ṭekrīthī nīce āvyũ} — The sheep came down from the hill.
  \item \emph{Recall:} say the Gujarati for to moo, then the Gujarati for to milk, then say \emph{gheṭũ} again
\end{itemize}

\begin{tcolorbox}[breakable,colback=teal!4,colframe=teal!35!black,title={Before you move on}]
What does \gu{ઘેટું} mean? (\textquotedblleft{}Sheep\textquotedblright{}.) Say it once more.
\end{tcolorbox}
\section[\texorpdfstring{\gu{ચરવું} (carvũ)}{ચરવું (carvũ)}]{\gu{ચરવું} (carvũ) — to graze}

\label{lesson:GU-C234-carvu}



\begin{quote}
Before the new one: say the Gujarati for to milk, then the Gujarati for sheep.
\end{quote}

\subsection*{You'll want to know: \gu{ચરવું}}
\textbf{\gu{ચરવું}} — \emph{carvũ} — \textquotedblleft{}to graze\textquotedblright{}.

\begin{quote}
\textbf{\gu{ઘેટાં} \gu{મેદાનમાં} \gu{ચર્યાં}.} — \emph{gheṭā̃ medānmā̃ caryā̃} — \textquotedblleft{}The sheep grazed in the meadow.\textquotedblright{}
\end{quote}

\begin{cousinweb}
\textbf{\gu{ચરવું}} (\emph{carvũ}) is a tadbhava. It comes from Sanskrit \emph{carati}, \textquotedblleft{}moves; grazes\textquotedblright{}, and CDIAL glosses the Gujarati verb \textquotedblleft{}to move about, graze\textquotedblright{} (CDIAL 4686).
\end{cousinweb}

\begin{grammarlens}[title={What you've built}]
One more word for everyday life at A2.
\end{grammarlens}

\subsection*{Your turn}
\begin{itemize}
\raggedright
  \item \emph{Say it:} \emph{carvũ}
  \item \emph{Say it:} \emph{carvũ}, once more
  \item \emph{Say it:} \emph{gheṭā̃ medānmā̃ caryā̃} — The sheep grazed in the meadow.
  \item \emph{Recall:} say the Gujarati for to milk, then the Gujarati for sheep, then say \emph{carvũ} again
\end{itemize}

\begin{tcolorbox}[breakable,colback=teal!4,colframe=teal!35!black,title={Before you move on}]
What does \gu{ચરવું} mean? (\textquotedblleft{}To graze\textquotedblright{}.) Say it once more.
\end{tcolorbox}
